%==============================================================
% BAB II  TINJAUAN PUSTAKA
%==============================================================

\chapter{TINJAUAN PUSTAKA}

%--------------------------------------------------------------
\section{Retinopathy of Prematurity dan Tahapannya}
%--------------------------------------------------------------

\textit{Retinopathy of Prematurity} (ROP) adalah kelainan pembuluh darah retina yang terjadi
pada bayi yang lahir sebelum waktunya, terutama yang lahir sebelum usia kehamilan 31 minggu
atau dengan berat badan lahir di bawah 1.500 gram. Kondisi ini terjadi karena pembuluh darah
retina yang belum sempat terbentuk secara lengkap di dalam kandungan kemudian tumbuh secara
tidak normal setelah kelahiran, sebagai respons terhadap kadar oksigen yang berubah dan
faktor pertumbuhan yang tidak seimbang \citep{Chiang2021}.

Berdasarkan Klasifikasi Internasional ROP (\textit{International Classification of
Retinopathy of Prematurity}, ICROP) yang diperbarui pada tahun 2021 \citep{Chiang2021},
keparahan ROP dibagi ke dalam lima \textit{stage}:

\begin{itemize}
    \item \textbf{Normal}: Pembuluh darah retina berkembang secara teratur dari area \textit{optic
    disc} ke arah pinggiran retina.
    \item \textbf{Stage 1}: Terdapat garis batas datar (\textit{demarcation line}) yang
    memisahkan retina yang sudah berpembuluh darah dari yang belum.
    \item \textbf{Stage 2}: Garis batas tersebut menebal dan terangkat membentuk tonjolan
    jaringan (\textit{ridge}).
    \item \textbf{Stage 3}: Terjadi pertumbuhan pembuluh darah baru ke arah \textit{vitreous} di
    atas \textit{ridge} (neovaskularisasi ekstraretina).
    \item \textbf{Stage 4--5}: Ablasio retina sebagian atau menyeluruh (tidak tercakup
    dalam penelitian ini).
\end{itemize}

Selain itu, dikenal juga kondisi \textit{plus disease} yang ditandai oleh dilatasi dan
\textit{tortuositas} pembuluh darah di kutub posterior retina, yang dapat menyertai \textit{stage}
manapun dan menandakan penyakit yang lebih aktif \citep{Chiang2021}. Perbedaan visual antar
\textit{stage} awal, khususnya antara \textit{Stage} 1 dan \textit{Stage} 2, sangat halus
sehingga membutuhkan keahlian khusus dan pengalaman klinis yang tinggi \citep{Vahidmoghadam2026}.
Ilustrasi setiap \textit{stage} ditunjukkan pada Gambar \ref{fig:rop_stages_diagram}.

\begin{figure}[H]
  \centering
  \includegraphics[width=0.85\linewidth]{images/rop_stages_diagram}
  \caption{Ilustrasi klasifikasi kategori kelas retina pada penyakit ROP}
  \label{fig:rop_stages_diagram}
\end{figure}

%--------------------------------------------------------------
\section{Dataset Citra Fundus ROP}
%--------------------------------------------------------------

Kemajuan dalam klasifikasi ROP berbasis komputer sangat bergantung pada ketersediaan dataset
citra yang berlabel dan berkualitas baik. \citet{Zhao2024} menerbitkan dataset citra fundus ROP
pertama yang bersifat publik dengan skala besar, yang terdiri dari 1.099 citra fundus bayi
prematur yang dikumpulkan dari rumah sakit di Tiongkok menggunakan kamera RetCam. Dataset ini
mencakup lima kategori: Normal (236 citra), Stage 1 (94 citra), Stage 2 (165 citra), Stage 3
(261 citra), dan bekas luka laser (343 citra). Setiap citra dilengkapi dengan label \textit{stage}
yang ditentukan oleh dokter spesialis mata berpengalaman. Dataset ini juga disertai dengan model
\textit{benchmark} berbasis ResNet50 pra-latih yang mencapai \textit{macro F1-score} sebesar
0,8281 pada pengaturan 5-kelas \citep{Zhao2024}.

%--------------------------------------------------------------
\section{Praproses dan Enhancement Citra Fundus}
%--------------------------------------------------------------

Kualitas citra fundus bayi umumnya lebih rendah dibanding citra fundus orang dewasa karena
berbagai faktor teknis dan kondisi pasien. Pencahayaan yang tidak merata, refleksi kornea yang
kuat, dan kontras yang rendah antara pembuluh darah dan jaringan sekitarnya merupakan tantangan
utama yang perlu ditangani sebelum proses analisis otomatis \citep{Rahim2024}.

\citet{Rahim2024} menunjukkan bahwa kombinasi koreksi pencahayaan adaptif dan CLAHE pada ruang
warna L*a*b* secara signifikan meningkatkan visibilitas struktur klinis penting seperti
\textit{demarcation line} dan \textit{ridge}. CLAHE bekerja dengan membagi citra menjadi
blok-blok kecil dan menerapkan pemerataan histogram secara terpisah pada setiap blok, dengan
pembatasan kliping (\textit{clip limit}) untuk mencegah penguatan berlebihan pada area homogen.
Penggunaan ruang warna L*a*b* memungkinkan \textit{enhancement} diterapkan hanya pada komponen
kecerahan (L*) tanpa mengubah informasi warna (a* dan b*) \citep{Rahim2024}.

Normalisasi latar belakang dengan filter Gaussian berskala besar juga merupakan teknik yang
efektif untuk mengurangi ketimpangan pencahayaan. Dengan memperkirakan distribusi pencahayaan
menggunakan filter Gaussian berukuran besar, kemudian membaginya dengan distribusi tersebut,
variasi kecerahan antar area citra dapat dikurangi secara substansial tanpa merusak informasi
detail lokal \citep{Almeida2024}.

%--------------------------------------------------------------
\section{Representasi Pembuluh Darah dan Ridge Retina}
%--------------------------------------------------------------

Representasi struktur pembuluh darah dan \textit{ridge} merupakan bagian penting dalam
analisis citra fundus untuk ROP. Pembuluh darah retina memiliki karakteristik geometris
yang khas, yaitu struktur tubular dengan intensitas yang lebih rendah dibanding jaringan
sekitarnya pada kanal hijau citra RGB, sedangkan \textit{ridge} dan \textit{demarcation line}
muncul sebagai struktur tepi yang membedakan stage awal ROP.

\subsection{Metode Berbasis Filter Spasial}

Pendekatan berbasis filter spasial memanfaatkan karakteristik geometris pembuluh darah untuk
membedakannya dari latar belakang. Filter Gabor multiskala dan multiorientasi terbukti efektif
dalam mendeteksi struktur tubular pada berbagai orientasi. Filter ini merespons secara kuat
terhadap tepi berarah yang merupakan ciri khas pembuluh darah retina \citep{Khandouzi2022}.

Filter berbasis matriks Hessian, seperti filter Frangi, Sato, dan Meijering, memanfaatkan
nilai eigen dari matriks Hessian orde dua untuk mengukur derajat keserupaan dengan pembuluh darah (\textit{vesselness}) pada
setiap piksel. Filter Meijering, yang dikembangkan oleh \citet{Meijering2004} untuk pelacakan
filamen neurit dalam citra fluoresensi, dirancang khusus untuk mendeteksi struktur filamen tipis
dan terbukti bekerja lebih baik pada pembuluh darah yang halus dibanding filter Frangi
\citep{Tian2021}. \citet{Almeida2024} menunjukkan bahwa kombinasi filter berbasis Hessian
dengan teknik \textit{top-hat} morfologis mampu menghasilkan peta \textit{vesselness} dengan
respons yang kuat pada struktur tubular tanpa memicu terlalu banyak deteksi palsu.

\subsection{Representasi Softmap Struktural}

Pada tugas klasifikasi, respons filter spasial tidak harus diubah menjadi masker biner.
Respons tersebut dapat dipertahankan sebagai \textit{softmap} kontinu agar jaringan saraf tetap
menerima informasi kekuatan respons lokal. Dalam penelitian ini, filter Gabor digunakan untuk
menonjolkan tekstur pembuluh darah, sedangkan respons Hessian berbasis Meijering digunakan
bersama \textit{oriented top-hat} untuk memperkuat struktur garis batas dan \textit{ridge}.
Pemilahan kanal seperti ini membuat representasi input tidak hanya membawa informasi vaskular,
tetapi juga struktur tepi yang relevan untuk membedakan Stage 1, Stage 2, dan Stage 3.

%--------------------------------------------------------------
\section{Klasifikasi ROP Berbasis \textit{Deep Learning}}
%--------------------------------------------------------------

Penggunaan jaringan saraf tiruan untuk analisis citra medis ROP telah berkembang pesat dalam
beberapa tahun terakhir. \citet{Zhang2021} melakukan tinjauan sistematis dan meta-analisis
terhadap 24 penelitian tentang penggunaan algoritma \textit{deep learning} untuk mendeteksi
ROP dari citra fundus, dan menemukan bahwa sebagian besar model mencapai akurasi di atas 90\%
untuk tugas biner (ROP vs. Normal), namun performa turun secara signifikan pada tugas
klasifikasi multi-\textit{stage}.

\citet{Tong2020} mengembangkan sistem berbasis CNN untuk identifikasi otomatis ROP dan mengujinya
pada dataset klinis dari rumah sakit di Tiongkok. Model mereka mencapai sensitivitas dan
spesifisitas yang tinggi untuk deteksi ROP membutuhkan terapi, namun tidak membahas secara
detail performa per-\textit{stage}.

\citet{Vahidmoghadam2026} mengusulkan pendekatan yang menggabungkan informasi citra dengan
masker pembuluh darah untuk mendeteksi \textit{plus disease} dan \textit{stage} awal ROP secara
bersamaan. Penelitian ini menunjukkan bahwa informasi vaskular yang disajikan sebagai masukan
tambahan dapat meningkatkan kemampuan model dalam membedakan \textit{stage} yang secara visual
sangat mirip.

\subsection{Jaringan Residual (\textit{Residual Network})}

Arsitektur jaringan saraf dengan \textit{residual connection} yang diperkenalkan oleh
\citet{He2016} merupakan dasar dari sebagian besar model \textit{deep learning} modern.
\textit{Residual connection} memungkinkan sinyal gradien mengalir langsung dari lapisan yang
lebih dalam ke lapisan yang lebih dangkal selama pelatihan, sehingga mengatasi masalah
\textit{vanishing gradient} dan memungkinkan pelatihan jaringan yang jauh lebih dalam. Untuk
dataset medis berukuran kecil, varian ResNet yang lebih kecil (\textit{tiny}) lebih sesuai
karena memiliki jumlah parameter yang lebih sedikit sehingga risiko \textit{overfitting} lebih
rendah.

\subsection{Penanganan Ketidakseimbangan Kelas}

Ketidakseimbangan jumlah sampel antar kelas merupakan masalah umum dalam dataset medis.
Pada dataset Zhao2024, perbandingan antara kelas terbanyak (bekas luka laser: 343 sampel) dan
kelas paling sedikit (Stage 1: 94 sampel) hampir mencapai 4:1. \citet{Lin2017} mengusulkan
\textit{focal loss} sebagai alternatif \textit{cross-entropy}---fungsi kerugian standar yang mengukur perbedaan antara distribusi probabilitas prediksi model dan label sebenarnya---yang secara otomatis mengurangi bobot pada sampel yang mudah diklasifikasikan sehingga pelatihan lebih berfokus pada sampel yang sulit, yang biasanya merupakan sampel dari kelas minoritas. Pendekatan berbasis
\textit{ensemble} seperti \textit{Random Forest} juga mendukung mekanisme pembobotan kelas secara
\textit{native} \citep{Breiman2001}.

Selain itu, pemberian bobot kelas berbanding terbalik dengan frekuensinya (\textit{inverse
frequency weighting}) merupakan teknik sederhana namun efektif untuk mengatasi
ketidakseimbangan kelas \citep{Lin2017}. Teknik ini memastikan bahwa kesalahan pada kelas
minoritas mendapat penalti yang lebih besar selama pelatihan.

\subsection{Validasi Silang yang Sadar Kelompok (\textit{Group-Aware Cross-Validation})}

Estimasi performa model yang andal membutuhkan protokol validasi yang tepat. Pada dataset
citra medis, satu pasien atau satu mata dapat menghasilkan banyak citra yang sangat mirip.
Apabila data dibagi secara acak pada tingkat citra, citra dari mata yang sama berpotensi masuk
ke sisi pelatihan sekaligus sisi pengujian, sehingga model dapat mengenali identitas pasien
alih-alih mempelajari fitur penyakit dan estimasi performa menjadi terlalu optimis.
\textit{Group-aware cross-validation} (\textit{GroupKFold}) mengatasi risiko ini dengan
memastikan seluruh citra dari kelompok yang sama selalu berada di sisi yang sama selama
validasi silang, sehingga performa yang dilaporkan mencerminkan kemampuan generalisasi model
terhadap pasien yang benar-benar belum pernah dilihat \citep{Zhao2024}.
